% HD-G34. GPT-6 Astra Pro /root/g34; inherited effort not exposed.
% 21 November 2026. Actual primitive rank-four source exclusion.
\section{HD-G34: closure of the primitive rank-four branch}
\label{sec:primitive-rank-four-finite-closure}

Retain the original natural Weil-sixfold source with exact
coefficient character field \(K\) and period rank nine, where
\[
 K=\mathbb Q(i)\quad\hbox{or}\quad K=\mathbb Q(\sqrt{-d}),
 \qquad d\in\{2,7,11,19,43,67,163\}.
\]
Suppose its actual finite
coefficient group has no nonidentity element of moving rank
one, two or three. Its normal subgroup generated by all actual
moving-rank-four elements cannot be absolutely irreducible and
complex primitive. This closes the primitive branch left by
HD-R18. The argument takes place in the original coefficient
representation and its actual normalizer.

\paragraph{1. Actual powers and conjugate generation.}
Write \(V=K^m\), \(G\subset GL(V)\), and retain from HD-R18
the characteristic absolutely irreducible quasi-simple subgroup
\(S\triangleleft N_4\triangleleft G\). Put
\(L=S/Z(S)\) and \(Q=G\cap K^*I\). The proved actual maps give
\[
 21\le m\le63,\qquad G/Q\hookrightarrow\operatorname{Aut}(L),
 \qquad |Z(S)|\mid\gcd(m,|\mu(K)|).                 \tag{G34.1}
\]
For the original nonidentity branch matrices, let \(c_j\)
be their moving ranks. The sphere Euler equation and the
dimension of the marked sphere give
\[
       \sum_{j=1}^b c_j=2m+6,\qquad b\ge12,
       \qquad c_j\ge4.                              \tag{G34.2}
\]

Take an actual moving-rank-four element \(x\). All four
nonunit eigenvalues have the same order \(a>1\). Indeed, two
different orders can be labelled \(a_1,a_2\) with
\(a_1\nmid a_2\). Raising \(x\) to the power \(a_2\) fixes at least one,
but not all, of its four moving eigenvectors. Its moving rank
is then one, two or three, contrary to the stated property of
the original \(G\). For any prime \(p\mid a\), the actual
power \(y=x^{a/p}\) has order \(p\), with the same fixed space.
The four moving roots of \(y\) are primitive \(p\)-th roots.
Each irreducible factor of \(\Phi_p\) over \(K\) has degree
\[
 e=[K(\zeta_p):K]
   =\frac{p-1}{[K\cap\mathbb Q(\zeta_p):\mathbb Q]},
 \qquad e\mid4.                                    \tag{G34.3}
\]
The denominator is one or two. The only prime possibilities
are \(p=2,3,5\): if the denominator is two, then
\((p-1)/2\in\{1,2,4\}\), and nine is not prime.
Our fields contain neither \(\zeta_3\) nor the real quadratic
subfield \(\mathbb Q(\sqrt5)\) of \(\mathbb Q(\zeta_5)\).
Consequently the exact moving spectra and traces are
\[
\begin{array}{c|c|c}
p&\hbox{moving spectrum}&\operatorname{tr}(y)\\\hline
2&-1,-1,-1,-1&m-8\\
3&\zeta_3,\zeta_3,\zeta_3^2,\zeta_3^2&m-6\\
5&\zeta_5,\zeta_5^2,\zeta_5^3,\zeta_5^4&m-5.
\end{array}                                           \tag{G34.4}
\]
No scalar has been adjoined to the coefficient group.
Its image \(\bar y\) in (G34.1) is nonidentity of order \(p\),
since a nonsingular scalar in degree at least21 cannot have
moving rank four.

Use Guralnick--Saxl's exact definition of \(\alpha(\bar y)\),
the least number of \(L\)-conjugates generating
\(\langle L,\bar y\rangle\). Lift the conjugating elements
to the original \(S\), and let \(J\) be generated by the
resulting actual conjugates of \(y\). Then \(S\subset QJ\),
and perfection gives \(S\subset[QJ,QJ]=J'\).
Their common fixed space is therefore zero. Intersecting the
fixed spaces of these actual moving-rank-four matrices proves
\[
                       m\le4\alpha(\bar y).             \tag{G34.5}
\]

\paragraph{2. The complete corrected finite filter.}
We use Hiss--Malle, \emph{Low-dimensional representations of
quasi-simple groups}, LMS J. Comput. Math.4 (2001),22--63,
Theorem1 p.22 and Section6 p.30, with their complete replacement
table in \emph{Corrigenda}, volume5 (2002),95--126. All degree
21--63 rows lie on printed pp.100--108; the correction instructions
are on pp.95--96. The original generic Table2 p.31 remains in
force. Characteristic zero is included in a row labelled
\(\ell\ne p_1,p_2,\ldots\). Centres divisible by three are
excluded by (G34.1). Other character-field restrictions can be
ignored in making the following larger list.

For clarity, the exact coarse degree lists and all arithmetic
comparisons are recorded in the finite certificate. The uniform
conjugate bounds used are Guralnick--Saxl, J. Algebra268
(2003),519--571: Lemma3.1 pp.527--529; Lemmas3.2--3.4
pp.529--533; Theorems4.1--4.4 pp.533--534; Propositions5.4,
5.5,5.7,5.8 pp.544--547; Lemma6.1 pp.547--548; and the
relevant sporadic cases of Lemma7.6 pp.553--554 and Table1
p.554. For the Tits group the proof of Proposition5.5 at
\(q=2\), pp.545--546, gives \(\alpha\le4\); for
\(\operatorname{Sz}(8)\), Proposition5.8 gives \(\alpha\le3\).
Every sporadic ordinary row of degree at most63 with a permitted
centre has \(\alpha\le5\), so (G34.5) excludes it.

The generic \(L_2(q)\) rows have \(\alpha\le4\), apart
from \(L_2(9)=A_6\), for which the alternating bound is five;
neither permits degree at least21. In the generic alternating
rows, \(A_n\), \(n\ge14\), the ordinary module in this range
is the actual rational augmentation module of degree \(n-1\).
Solving its finite intertwiner equations over \(K\) identifies
the original restriction with that module: a nonzero solution
exists because the complex solution space has dimension one,
and it is invertible by irreducibility. An actual three-cycle
in \(S\) then has moving rank two, which is already forbidden.

Applying (G34.5) to all the other corrected rows leaves precisely
the following larger list:
\[
\begin{array}{c|l|c}
L&\hbox{possible degrees}&\hbox{uniform bound for }\alpha\\\hline
A_7&21&6\\
A_8&21,24,28&7\\
A_9&21,27,28&8\\
A_{10}&35,36&9\\
A_{12},A_{13}&32\ ;\ 32&11\ ;\ 12\\
U_4(2),U_4(3),U_6(2)&24\ ;\ 21\ ;\ 22&6\\
\operatorname{Sp}_6(2)&21,27&7\\
\operatorname{Sp}_8(2)&35&9\\
O_8^+(2)&28&8\\
{}^3D_4(2)&26&7.
\end{array}                                             \tag{G34.6}
\]
The degree24 alternating row is \(2.A_8\), and the degree32
rows are \(2.A_{12}\) and \(2.A_{13}\). The nonalternating
rows in (G34.6) are the simple groups displayed, not covers:
this is read from the same corrected degree rows. In particular,
the degree21 row of \(U_4(3)\) and the degree22 row of \(U_6(2)\)
are on p.101 and both have indicator \(+\). All the other
nonalternating rows in (G34.6) also have indicator \(+\):
\(U_4(2),24\) and \(\operatorname{Sp}_6(2),21\) on p.101;
\(\operatorname{Sp}_6(2),27\), \(O_8^+(2),28\),
and \({}^3D_4(2),26\) on p.102; and
\(\operatorname{Sp}_8(2),35\) on p.103.
Thus each of these actual restricted representations is self-dual.
The table indicators are used with their stated meaning in
Hiss--Malle Section5 pp.28--29, namely the existence of a
nondegenerate invariant symmetric bilinear form over \(\mathbb C\).
The descent of that form to the original \(K\)-representation is
proved next; it is not inferred from a different coefficient model.

\paragraph{3. The actual bilinear-form and branch-budget obstruction.}
Suppose the absolutely irreducible restriction to \(S\) is
self-dual. The equations
\[
                 \rho(s)^t B\rho(s)=B\quad(s\in S)       \tag{G34.7}
\]
are linear over \(K\). Their complex solution space is a line,
and every nonzero solution is nonsingular by irreducibility.
The coefficient matrix has the same rank over \(K\) and over
\(\mathbb C\), so there is such a nonsingular \(B\) over \(K\).
It can be symmetric or alternating; the following argument uses
only its nondegeneracy. Since the original \(G\) normalizes
the original \(S\), uniqueness in (G34.7) gives an actual
character
\[
 \rho(g)^tB\rho(g)=c(g)B,
 \quad c:G\longrightarrow\mu(K),\qquad
 B\rho(g)=c(g)\rho(g)^{-t}B.                            \tag{G34.8}
\]
The last formula is the literal isomorphism
\(V\to V^*\otimes c\); multiplicativity follows by applying
the first formula to a product, and finite order makes each
\(c(g)\) a root of unity in \(K\).

If \(c(g)\ne1\), its fixed space \(F\) is totally isotropic
for \(B\), since \(B(u,v)=c(g)B(u,v)\) for \(u,v\in F\).
The inclusion \(F\subset F^\perp\) and nondegeneracy give
\(\dim F\le m/2\). Thus that original matrix has moving rank
at least \(\lceil m/2\rceil\).

The character \(c\) is not trivial on all inertia generators.
Otherwise it is trivial on \(G\), and (G34.8) makes every
\(\rho(g)\) similar to its inverse transpose. Its trace then
equals its inverse trace, which is its complex-conjugate trace
because \(g\) has finite order. All coefficient traces would
belong to \(K\cap\mathbb R=\mathbb Q\), contradicting the
exact character field \(K\). The product-one relation on the
inertia generators implies that at least two of their
\(c\)-values differ from one. Combining their moving ranks
with (G34.2) gives the exact numerical obstruction
\[
 2m+6=\sum c_j
       \ge4(b-2)+2\lceil m/2\rceil
       \ge40+2\lceil m/2\rceil.                         \tag{G34.9}
\]
For every \(21\le m\le33\) the last quantity is strictly
greater than \(2m+6\). This excludes all the nonalternating
rows of (G34.6) except \((\operatorname{Sp}_8(2),35)\).

\paragraph{4. The alternating rows and their literal normalizer lifts.}
For every alternating row of (G34.6), the degree exceeds
\(4\lfloor n/2\rfloor\). Guralnick--Saxl Lemma6.1 and
(G34.5) therefore force \(\bar y\) to be a transposition in
\(\operatorname{Aut}(A_n)=S_n\). In particular \(y\) is
the actual order-two power in (G34.4).
Choose original \(S\)-conjugates \(T_i\) of \(y\) with images
\((i,i+1)\), for \(1\le i<n\). For any two of these matrices,
their fixed spaces have intersection dimension at least
\(m-8>0\). The elements \((T_iT_j)^2\), \(|i-j|>1\),
and \((T_iT_{i+1})^3\) lie in the actual scalar kernel \(Q\)
and fix that intersection pointwise. They are therefore one.
Together with \(T_i^2=1\), these are the exact Coxeter relations,
giving the actual injective homomorphism
\[
 \Phi:S_n\longrightarrow G,
 \qquad (i,i+1)\longmapsto T_i,
 \qquad \overline\Phi=\operatorname{id}_{S_n}.             \tag{G34.10}
\]
As in paragraph1, \(S\subset Q\Phi(S_n)\); perfection gives
\(S\subset\Phi(A_n)\). Surjectivity onto \(A_n\) and the
trivial scalar kernel of \(\Phi(A_n)\) yield
\(S=\Phi(A_n)\). Hence all double-cover rows disappear.

The actual \(S_n\)-module is irreducible because its restriction
to \(S\) is. The ordinary Young construction realizes every
irreducible \(S_n\)-module over \(\mathbb Q\), as the image
of right multiplication by its row--column symmetrizer on the
left regular module \(\mathbb Q[S_n]\).
For exact conventions and the primary proof, see Etingof et al.,
\emph{Introduction to representation theory}, Student
Mathematical Library59 (2011), the author version dated
10 January2011, Theorem4.36 and complete proof pp.58--60,
and Theorem4.53 and complete proof pp.63--64. With unnormalized
row--column element \(a_\lambda\), that projector is right
multiplication by \((\dim V_\lambda/n!)a_\lambda\), since
\(a_\lambda^2=(n!/\dim V_\lambda)a_\lambda\).
Averaging any rational positive definite form over \(S_n\)
gives a nondegenerate invariant rational symmetric form on that
actual image. The finite intertwiner equations over \(K\)
transport the representation and form to the original \(\Phi(S_n)\);
their nonzero solution is invertible by irreducibility. Its
restriction supplies the actual self-duality of \(S\), including
exclusion of any incompatible nonreal alternating character.
Thus (G34.9) also excludes all remaining alternating rows of
degree at most33.

Only \(n=10\), \(m=35,36\), remain. We record explicitly the
class-sum fact used here and in paragraph5. For a conjugacy class
\(\mathcal C\) of an actual finite group acting irreducibly on
\(V\), the central element \(z=\sum_{s\in\mathcal C}s\)
acts by the scalar
\[
           \kappa=\frac{|\mathcal C|\operatorname{tr}(s)}m.
                                                               \tag{G34.11}
\]
Left multiplication by \(z\) on the integral free module
\(\mathbb Z[\langle\mathcal C\rangle]\), or on the full
integral group ring, is an integral matrix. Its monic
characteristic polynomial annihilates \(z\), since left
multiplication is faithful. Applying the original representation
shows that \(\kappa\) is an algebraic integer. In particular a
rational \(\kappa\) must be an integer. This is an actual
integral group-algebra map, not a character-table assumption.

In the actual \(S_{10}\) of (G34.10), the45 transpositions are
one class. For \(m=35\) their required trace27 would give
\(\kappa=45\cdot27/35=243/7\), which is not an integer.
For \(m=36\), the hook formula has exactly two partitions of
ten of degree36, namely \((8,1,1)\) and \((3,1,1,1,1,1,1,1)\).
They are the following actual rational modules and their maps.
Set \(W=\ker(\mathbb Q^{10}\xrightarrow{\sum}\mathbb Q)\).
Then \(\Lambda^2W\) has dimension36. A transposition acts on
\(W\) with eight eigenvalues one and one eigenvalue minus one,
so its exterior-square trace is \(28-8=20\); its sign twist
has trace \(-20\). The character of \(\Lambda^2W\) at a
permutation with \(f\) fixed labels and \(t\) two-cycles is
\(((f-1)^2-(f+2t-1))/2\). Summing its square over the42 cycle
partitions of ten, with class weights
\(10!/\prod_j j^{a_j}a_j!\), gives norm one, and its scalar
product with its sign twist is zero. These exact finite sums
and the complete hook enumeration are recorded in the certificate.
The two actual rational models therefore exhaust the two
degree36 possibilities. Solving
\(C\rho(T_i)=R((i,i+1))C\) gives the actual invertible
\(K\)-intertwiner, so the original trace is \(20\) or \(-20\).
Neither equals the trace28 required by (G34.4).

\paragraph{5. The actual symplectic class sum.}
The remaining core has \(L=\operatorname{Sp}_8(2)\) and
\(m=35\). Its centre is trivial by (G34.1), so the original
\(S\) is itself isomorphic to that simple group.
Guralnick--Saxl Theorem4.3, pp.533--534, bounds
\(\alpha(x)\le7\) for every nontransvection in its full
automorphism group; the transvections have \(\alpha=9\).
Since \(35>4\cdot7\), (G34.5) forces \(\bar y\) to be an
inner transvection, of order two. Choose the unique \(s\in S\)
inducing that inner automorphism. Then \(s^2=1\), and the
actual centralizer of \(S\) in \(G\) is \(Q\) by absolute
irreducibility. Thus
\[
                     y=q s,\qquad q\in Q,\qquad q^2=1,
                 \qquad\operatorname{tr}(s)=\pm27.       \tag{G34.12}
\]
No assertion about an extension of a stored simple-group
character is needed.

Here is the exact class-size calculation in the standard
finite symplectic group. On \(\mathbb F_2^8\) with its
nondegenerate alternating form \(\langle\ ,\ \rangle\),
its transvections are precisely
\[
                  t_v(w)=w+\langle w,v\rangle v,
                       \qquad v\ne0.                    \tag{G34.13}
\]
Each is a nonidentity symplectic involution. Distinct nonzero
\(v\) give distinct rank-one image lines and hence distinct
matrices, since a line over \(\mathbb F_2\) has one nonzero
vector. A symplectic basis extending any chosen nonzero vector
shows transitivity on these vectors, and
\(a t_v a^{-1}=t_{av}\). Thus they form one conjugacy class
of size \(2^8-1=255\). Transporting (G34.13) through the
actual abstract isomorphism \(S\to\operatorname{Sp}_8(2)\)
identifies the class of the original \(s\); all class sums
below still act through its original \(K\)-matrices.
By (G34.11) this class sum has scalar
\[
                    \kappa=\frac{255(\pm27)}{35}
                            =\pm\frac{1377}{7}.          \tag{G34.14}
\]
It is rational and not integral, contradicting its integral
group-algebra characteristic polynomial. This excludes the
last row.

\paragraph{6. Scope and the original geometric factor.}
All contradictions above concern the original \(G\), its
actual inertia matrices and their actual powers or conjugates.
The bilinear map (G34.8), the Coxeter homomorphism (G34.10),
the finite \(K\)-intertwiners and the integral class-sum map
were specified explicitly. None removes an original normalizer
or replaces its coefficient character. A displayed intertwiner
acts componentwise on the usual parabolic cocycle complex,
with inverse componentwise \(C^{-1}\), and hence preserves its
kernel modulo coboundaries and its period rank. Equivalently,
one can keep the original complex throughout and use that
intertwiner only to evaluate the indicated original trace.
The covering curve and its actual group-algebra primitive
Jacobian factor are then literally the same objects, with
identity maps; no new polarized-isogeny comparison is used in
this exclusion. Together with HD-R18 this proves the stated
closure of the entire absolutely irreducible primitive
moving-rank-four branch.
